% Options for packages loaded elsewhere
\PassOptionsToPackage{unicode}{hyperref}
\PassOptionsToPackage{hyphens}{url}
\documentclass[
  10pt,
]{article}
\usepackage{xcolor}
\usepackage[margin=0.78in]{geometry}
\usepackage{amsmath,amssymb}
\setcounter{secnumdepth}{5}
\usepackage{iftex}
\ifPDFTeX
  \usepackage[T1]{fontenc}
  \usepackage[utf8]{inputenc}
  \usepackage{textcomp} % provide euro and other symbols
\else % if luatex or xetex
  \usepackage{unicode-math} % this also loads fontspec
  \defaultfontfeatures{Scale=MatchLowercase}
  \defaultfontfeatures[\rmfamily]{Ligatures=TeX,Scale=1}
\fi
\usepackage{lmodern}
\ifPDFTeX\else
  % xetex/luatex font selection
\fi
% Use upquote if available, for straight quotes in verbatim environments
\IfFileExists{upquote.sty}{\usepackage{upquote}}{}
\IfFileExists{microtype.sty}{% use microtype if available
  \usepackage[]{microtype}
  \UseMicrotypeSet[protrusion]{basicmath} % disable protrusion for tt fonts
}{}
\makeatletter
\@ifundefined{KOMAClassName}{% if non-KOMA class
  \IfFileExists{parskip.sty}{%
    \usepackage{parskip}
  }{% else
    \setlength{\parindent}{0pt}
    \setlength{\parskip}{6pt plus 2pt minus 1pt}}
}{% if KOMA class
  \KOMAoptions{parskip=half}}
\makeatother
\usepackage{graphicx}
\makeatletter
\newsavebox\pandoc@box
\newcommand*\pandocbounded[1]{% scales image to fit in text height/width
  \sbox\pandoc@box{#1}%
  \Gscale@div\@tempa{\textheight}{\dimexpr\ht\pandoc@box+\dp\pandoc@box\relax}%
  \Gscale@div\@tempb{\linewidth}{\wd\pandoc@box}%
  \ifdim\@tempb\p@<\@tempa\p@\let\@tempa\@tempb\fi% select the smaller of both
  \ifdim\@tempa\p@<\p@\scalebox{\@tempa}{\usebox\pandoc@box}%
  \else\usebox{\pandoc@box}%
  \fi%
}
% Set default figure placement to htbp
\def\fps@figure{htbp}
\makeatother
% definitions for citeproc citations
\NewDocumentCommand\citeproctext{}{}
\NewDocumentCommand\citeproc{mm}{%
  \begingroup\def\citeproctext{#2}\cite{#1}\endgroup}
\makeatletter
 % allow citations to break across lines
 \let\@cite@ofmt\@firstofone
 % avoid brackets around text for \cite:
 \def\@biblabel#1{}
 \def\@cite#1#2{{#1\if@tempswa , #2\fi}}
\makeatother
\newlength{\cslhangindent}
\setlength{\cslhangindent}{1.5em}
\newlength{\csllabelwidth}
\setlength{\csllabelwidth}{3em}
\newenvironment{CSLReferences}[2] % #1 hanging-indent, #2 entry-spacing
 {\begin{list}{}{%
  \setlength{\itemindent}{0pt}
  \setlength{\leftmargin}{0pt}
  \setlength{\parsep}{0pt}
  % turn on hanging indent if param 1 is 1
  \ifodd #1
   \setlength{\leftmargin}{\cslhangindent}
   \setlength{\itemindent}{-1\cslhangindent}
  \fi
  % set entry spacing
  \setlength{\itemsep}{#2\baselineskip}}}
 {\end{list}}
\usepackage{calc}
\newcommand{\CSLBlock}[1]{\hfill\break\parbox[t]{\linewidth}{\strut\ignorespaces#1\strut}}
\newcommand{\CSLLeftMargin}[1]{\parbox[t]{\csllabelwidth}{\strut#1\strut}}
\newcommand{\CSLRightInline}[1]{\parbox[t]{\linewidth - \csllabelwidth}{\strut#1\strut}}
\newcommand{\CSLIndent}[1]{\hspace{\cslhangindent}#1}
\setlength{\emergencystretch}{3em} % prevent overfull lines
\providecommand{\tightlist}{%
  \setlength{\itemsep}{0pt}\setlength{\parskip}{0pt}}
\usepackage{booktabs}
\usepackage{longtable}
\usepackage{array}
\usepackage{float}
\usepackage{graphicx}
\usepackage{xcolor}
\definecolor{GroupBlue}{HTML}{2878B5}
\newcommand{\includefigure}[3]{%
  \begin{figure}[H]
    \centering
    \IfFileExists{#1}{%
      \includegraphics[width=0.96\linewidth,height=0.70\textheight,keepaspectratio]{#1}%
    }{%
      \fbox{\parbox[c][0.22\textheight][c]{0.86\linewidth}{\centering Missing figure: \texttt{\detokenize{#1}}}}%
    }
    \caption{#2}
    \label{#3}
  \end{figure}%
}
\newcommand{\tightinputtable}[1]{%
  \begingroup\small\setlength{\tabcolsep}{4pt}\input{#1}\endgroup%
}
\usepackage{bookmark}
\IfFileExists{xurl.sty}{\usepackage{xurl}}{} % add URL line breaks if available
\urlstyle{same}
% fallback for those not using the hyperref driver hyperxmp:
\makeatletter
\@ifundefined{xmpquote}{\newcommand{\xmpquote}[1]{#1}}{}
\makeatother
\hypersetup{
  pdftitle={STAT452 Final Project: Wine Quality and Experimental Design},
  pdfauthor={Trần Lê Anh Tuấn (24125107); Nguyễn Bảo Minh Triết (24125047); Nguyễn Đình Thiên Lộc (24125093); Nguyễn Hồng Tấn Tài (24125078); Lê Minh Thuận (24125105)},
  hidelinks,
  pdfcreator={LaTeX via pandoc}}

\title{STAT452 Final Project: Wine Quality and Experimental Design}
\author{Trần Lê Anh Tuấn (24125107) \and Nguyễn Bảo Minh Triết
(24125047) \and Nguyễn Đình Thiên Lộc (24125093) \and Nguyễn Hồng Tấn
Tài (24125078) \and Lê Minh Thuận (24125105)}
\date{Academic Year 2025--2026}

\begin{document}
\maketitle

{
\setcounter{tocdepth}{2}
\tableofcontents
}
\begin{abstract}
This report currently completes Part 1 of the Group 8 final project. We predict
red Vinho Verde quality from 11 physicochemical measurements with a locked,
leakage-controlled comparison of OLS, Ridge, Lasso, and Elastic Net. Part 2,
the separate experimental-design extension, remains explicitly pending.
\end{abstract}

\section{Part 1: Wine Quality (Red)}\label{part-1-wine-quality-red}

\section{Exploratory Data Analysis}\label{exploratory-data-analysis}

\subsection{Dataset, variables, and prediction
question}\label{dataset-variables-and-prediction-question}

The UCI Wine Quality data describe red Portuguese Vinho Verde samples
using 11 continuous physicochemical measurements and a sensory score
supplied as the median of at least three expert ratings
(\citeproc{ref-Cortez2009}{Cortez et al. 2009a},
\citeproc{ref-UCIWine}{2009b}). The response is \texttt{quality}; the
predictors cover acidity, sugar, chlorides, sulfur dioxide, density, pH,
sulphates, and alcohol. All predictors are numeric, and the detailed
data dictionary appears in Appendix A.

We model quality numerically because regression preserves its ordering
and avoids an arbitrary good/poor threshold. This assumes that a
one-point change has roughly comparable meaning across the observed 3--8
range. Gaussian regression can produce fractional expected scores, so
the interpretation is prediction of expected expert score rather than
exact ordinal classification.

After the integrity checks in Section 2, the unique records were divided
by a fixed quality-stratified 80/20 split. Every plot below uses only
the 1088 training observations; the 271 holdout outcomes were not
inspected.

\subsection{Response and predictor
distributions}\label{response-and-predictor-distributions}

\includefigure{../output/figures/fig_p2_quality_distribution.pdf}{Training-set distribution of the sensory response. Scores 5 and 6 dominate, while the extremes contain few observations.}{fig:p2-quality}

The response is concentrated at 5 and 6. A mean-like prediction is
therefore a strong baseline, and error estimates for rare scores 3, 4,
and 8 will be less precise. This imbalance also foreshadows regression
toward the center.

\includefigure{../output/figures/fig_p2_predictor_distributions.pdf}{Training-set predictor distributions before transformation. Several chemistry measures are strongly right-skewed.}{fig:p2-distributions}

Residual sugar, chlorides, sulfur-dioxide measures, and sulphates have
long right tails. These distributions motivate the transformations
documented in Section 2; EDA itself is shown on the original measurement
scale.

\subsection{Pairwise relationships and correlation
structure}\label{pairwise-relationships-and-correlation-structure}

\includefigure{../output/figures/fig_p2_correlation_heatmap.pdf}{Training-set Pearson correlations. Blocks among acidity, density, pH, and sulfur-dioxide measures motivate coefficient stabilization.}{fig:p2-correlation}

The strongest marginal response signals are alcohol (\(r=0.47\));
volatile acidity (\(r=-0.37\)); sulphates (\(r=0.23\)); citric acid
(\(r=0.22\)). The heatmap also shows correlated predictor blocks,
especially among acidity, density, pH, and the two sulfur-dioxide
measurements. That structure makes OLS coefficients potentially unstable
and directly motivates Ridge, Lasso, and Elastic Net.

\includefigure{../output/figures/fig_p2_pairwise_quality.pdf}{The four strongest marginal predictor--quality relationships in training. Jitter reveals the discrete response; lines are descriptive OLS trends.}{fig:p2-pairwise}

Alcohol is positively associated with quality and volatile acidity
negatively associated, but every score overlaps substantially on every
predictor. No single physicochemical measurement is sufficient. These
are predictive associations, not causal effects, because grape variety,
brand, vintage, price, and production conditions are absent.

\section{Data Cleaning}\label{data-cleaning}

\subsection{Missingness, duplicates, and the locked
split}\label{missingness-duplicates-and-the-locked-split}

The raw archive contains 1599 rows, 11 numeric predictors, no missing
cells, and no non-finite values. It also contains 240 exact duplicate
records. Because the archive provides no sample identifier, an exact
repeat could be either a copied row or a genuinely repeated batch. We
conservatively keep the first copy and remove the rest \textbf{before}
splitting, preventing identical profiles and labels from appearing in
both training and holdout data. The limitation is that this choice may
underweight genuinely repeated batches.

\tightinputtable{../output/tables/tab_p1_data_audit.tex}

The fixed seed 4520803 produces a quality-stratified 80/20 partition
after deduplication. Seed 4520804 then assigns five stratified folds
inside training. The shared fold IDs are reused by every model.

\tightinputtable{../output/tables/tab_p1_split_summary.tex}

\subsection{Outliers and
transformations}\label{outliers-and-transformations}

Tukey's 1.5-IQR rule flags 277 training rows (25.5\%). Deleting all of
them would discard a large and scientifically plausible part of the
target population; the dataset metadata also warns that rare wines are
expected. We therefore retain every row, winsorize each predictor at its
training 1st and 99th percentiles to limit leverage, and never clip the
response.

\includefigure{../output/figures/fig_p2_outlier_audit.pdf}{Training observations flagged by the univariate Tukey-IQR rule. Flags are diagnosed and retained rather than deleted wholesale.}{fig:p2-outliers}

The training skewness rule selects residual sugar, chlorides, free
sulfur dioxide, total sulfur dioxide, sulphates for \texttt{log1p}
transformation. Median imputation is defined as a reproducibility
safeguard but is a no-op for this dataset. After winsorization and
optional \texttt{log1p}, all predictors are centered and scaled so
penalty magnitudes are comparable.

\subsection{Leakage control}\label{leakage-control}

All medians, winsorization limits, transformation choices, means, and
standard deviations are learned without holdout outcomes. During
cross-validation, those quantities are refit inside each analysis fold
and applied to its validation fold. The final recipe is fit once to all
training predictors and then applied unchanged to the predictor-only
holdout file. This design prevents preprocessing information from
leaking backward into feature selection or tuning.

\section{Feature Selection}\label{feature-selection}

\subsection{Screening without deleting correlated
information}\label{screening-without-deleting-correlated-information}

We use three training-only views of relevance: marginal correlations,
variance inflation factors (VIFs), and an AIC stepwise sensitivity
analysis. The largest VIF is 7.90 for fixed acidity. This is material
multicollinearity but not evidence that the feature is useless: deleting
one correlated chemistry measure can make the retained coefficient
depend arbitrarily on that choice. Therefore all predictors enter Ridge,
while sparse selection is delegated to Lasso's embedded penalty.

\includefigure{../output/figures/fig_p3_feature_screening.pdf}{VIFs and full-OLS coefficients after the training-derived transformations. Coefficients are comparable because predictors are standardized.}{fig:p3-screening}

The AIC sensitivity model retains 7 of 11 predictors: alcohol, volatile
acidity, sulphates, chlorides, ph, total sulfur dioxide, free sulfur
dioxide. We do not promote this data-dependent subset directly to the
holdout comparison; it is a diagnostic against which to compare Lasso.

\subsection{Mean and OLS baselines}\label{mean-and-ols-baselines}

\tightinputtable{../output/tables/tab_p3_baseline_performance.tex}

The mean-only model establishes the no-chemistry benchmark. OLS lowers
cross-validated error, but correlated predictors allow coefficient
variance. Regularization deliberately accepts some training bias to
reduce that variance.

\subsection{Embedded Lasso selection}\label{embedded-lasso-selection}

At the one-standard-error penalty, Lasso retains 7 predictors: alcohol,
volatile acidity, sulphates, chlorides, ph, total sulfur dioxide, fixed
acidity. This is the required principled feature-selection result. The
one-SE rule is used for interpretation because it favors a stable sparse
description; the minimum-CV penalty remains the predeclared predictive
fit.

\tightinputtable{../output/tables/tab_p4_lasso_selection.tex}

\section{Modelling with
Regularization}\label{modelling-with-regularization}

For standardized transformed predictors, all regularized candidates
minimize

\[
\frac{1}{2n}\lVert \mathbf y-\beta_0-\mathbf X\boldsymbol\beta\rVert_2^2
+\lambda\left\{\alpha\lVert\boldsymbol\beta\rVert_1
+\frac{1-\alpha}{2}\lVert\boldsymbol\beta\rVert_2^2\right\}.
\]

Ridge uses \(\alpha=0\) (\citeproc{ref-Hoerl1970}{Hoerl and Kennard
1970}), Lasso \(\alpha=1\) (\citeproc{ref-Tibshirani1996}{Tibshirani
1996}), and Elastic Net searches \(\alpha\in\{0.25,0.50,0.75\}\)
(\citeproc{ref-Zou2005}{Zou and Hastie 2005}). For each alpha, 120
decreasing \(\lambda\) values from 100 to \(10^{-4}\) use identical
folds. Coordinate-descent fits are produced with \texttt{glmnet}
(\citeproc{ref-Friedman2010}{Friedman et al. 2010}).

\subsection{Cross-validation and
shrinkage}\label{cross-validation-and-shrinkage}

\includefigure{../output/figures/fig_p4_cv_curves.pdf}{Five-fold training CV curves with fold-specific preprocessing. Dashed lines minimize CV error; dotted lines apply the one-standard-error rule.}{fig:p4-cv}

Elastic Net selected \(\alpha=0.75\). The minimum-error and one-SE
penalties, CV errors, and effective model sizes are reported below.
Ridge keeps all coefficients but continuously shrinks them; Lasso and
Elastic Net can set coefficients exactly to zero.

\tightinputtable{../output/tables/tab_p4_regularized_display.tex}

\includefigure{../output/figures/fig_p4_coefficient_paths.pdf}{Coefficient paths as regularization changes. The vertical line in each panel is the minimum-CV penalty used for prediction.}{fig:p4-paths}

As \(\lambda\) grows, paths move toward zero. Lasso reaches exact zeros
earliest; Ridge shares weight among correlated measurements; Elastic Net
interpolates. This is the expected bias--variance mechanism rather than
merely a software choice.

\subsection{Pre-holdout decision}\label{pre-holdout-decision}

Before reconstructing any holdout outcome, we locked \textbf{Lasso},
whose five-fold CV RMSE was 0.6712. This written lock prevents choosing
a preferred model after seeing test performance. All predeclared models
are nevertheless evaluated to make the comparison transparent.

\section{Evaluation}\label{evaluation}

\subsection{Fair model comparison}\label{fair-model-comparison}

\tightinputtable{../output/tables/tab_p5_holdout_display.tex}

On the same 271 rows, the pre-locked Lasso obtained RMSE 0.614, MAE
0.482, and \(R^2\) 0.439. The observed lowest holdout RMSE belongs to
\textbf{OLS} (0.613). We retain the CV lock as the confirmatory choice
even if this ranking differs: selecting by the holdout would make its
error optimistic. The bootstrap intervals are wide enough that small
ranking differences should not be overinterpreted.

\includefigure{../output/figures/fig_p5_actual_vs_predicted.pdf}{Held-out predictions from the four fitted regression models. The diagonal denotes perfect prediction.}{fig:p5-predictions}

Predictions shrink toward the common scores 5--6, which lowers average
error but systematically misses rare extremes. This pattern follows from
response imbalance and squared-error fitting, not simply from
regularization.

\subsection{Diagnostics and bias--variance
trade-off}\label{diagnostics-and-biasvariance-trade-off}

\includefigure{../output/figures/fig_p5_locked_residual_diagnostics.pdf}{Residual checks for the pre-locked model on the untouched holdout set. These diagnostics describe generalization and were not used to retune the model.}{fig:p5-residuals}

OLS has the least shrinkage and therefore the lowest regularization bias
but the greatest sensitivity to correlated predictors. Ridge reduces
coefficient variance without selection; Lasso trades additional bias for
sparsity; Elastic Net stabilizes groups of correlated features while
retaining some selection. The relevant empirical comparison is the gap
among training, CV, and holdout RMSE in Table
\ref{tab:p5-holdout-display}. The regularized candidates have similar
held-out accuracy, so stability and interpretability matter more than a
small decimal-place difference.

\subsection{Limitations}\label{limitations}

\begin{enumerate}
\def\labelenumi{\arabic{enumi}.}
\tightlist
\item
  Quality is ordinal and discrete; Gaussian regression treats adjacent
  scores as equally spaced and permits fractional predictions.
\item
  Rare scores 3, 4, and 8 yield imprecise tail performance, while the
  bootstrap quantifies only sampling variation in this holdout.
\item
  Exact deduplication prevents leakage but may remove genuinely repeated
  batches because the archive has no sample identifier.
\item
  Physicochemical measurements omit brand, vintage, grape variety,
  price, and production process; predictive associations are not causal
  effects.
\item
  Winsorization limits leverage and log transforms reduce skew, but
  alternative robust losses or ordinal models could change tail
  behavior.
\end{enumerate}

\subsection{Conclusion}\label{conclusion}

The full workflow meets the Part 1 objective: documented EDA and
cleaning, principled selection, a baseline plus three regularized
models, shared-fold CV, coefficient-path interpretation, and a single
locked holdout evaluation. Lasso is the confirmatory recommendation for
expected-score prediction under this design. Alcohol and volatile
acidity provide the strongest marginal signals, but correlated chemistry
measurements and irreducible sensory variation limit accuracy. A future
extension should compare ordinal regression or robust nonlinear models
under the same locked resampling protocol.

\clearpage

\section{Part 2: Experimental-design
extension}\label{part-2-experimental-design-extension}

TODO: Select and justify the separate dataset; state hypotheses;
complete the two-factor ANOVA or \(2^k\) factorial analysis; check
assumptions; report interactions and post-hoc comparisons where
justified; and interpret practical significance and limitations.

\section{Contributions}\label{contributions}

TODO: Replace the workstream ownership plan with a statement of the work
each member actually completed and reviewed.

\clearpage
\appendix

\section{Detailed data documentation}\label{detailed-data-documentation}

\tightinputtable{../output/tables/tab_p1_data_dictionary.tex}

\begingroup\scriptsize\setlength{\tabcolsep}{2.5pt}
\begin{table}

\caption{\label{tab:p2-descriptive-statistics}Training-set descriptive statistics.}
\centering
\begin{tabular}[t]{lrrrrrrrrr}
\toprule
Variable & N & Mean & SD & Min & Q1 & Median & Q3 & Max & Skewness\\
\midrule
fixed\_acidity & 1088 & 8.316 & 1.744 & 4.600 & 7.100 & 7.900 & 9.225 & 15.900 & 0.925\\
volatile\_acidity & 1088 & 0.531 & 0.185 & 0.160 & 0.390 & 0.520 & 0.640 & 1.580 & 0.808\\
citric\_acid & 1088 & 0.273 & 0.195 & 0.000 & 0.090 & 0.260 & 0.430 & 0.790 & 0.263\\
residual\_sugar & 1088 & 2.521 & 1.331 & 0.900 & 1.900 & 2.200 & 2.600 & 15.500 & 4.677\\
chlorides & 1088 & 0.088 & 0.048 & 0.012 & 0.070 & 0.079 & 0.090 & 0.611 & 5.421\\
\addlinespace
free\_sulfur\_dioxide & 1088 & 15.830 & 10.345 & 1.000 & 7.000 & 14.000 & 21.000 & 68.000 & 1.192\\
total\_sulfur\_dioxide & 1088 & 46.691 & 33.378 & 6.000 & 22.000 & 38.000 & 63.000 & 289.000 & 1.610\\
density & 1088 & 0.997 & 0.002 & 0.990 & 0.996 & 0.997 & 0.998 & 1.004 & 0.013\\
ph & 1088 & 3.310 & 0.156 & 2.860 & 3.210 & 3.310 & 3.400 & 4.010 & 0.258\\
sulphates & 1088 & 0.660 & 0.172 & 0.330 & 0.550 & 0.620 & 0.730 & 1.980 & 2.338\\
\addlinespace
alcohol & 1088 & 10.456 & 1.101 & 8.400 & 9.500 & 10.200 & 11.200 & 14.900 & 0.854\\
quality & 1088 & 5.625 & 0.825 & 3.000 & 5.000 & 6.000 & 6.000 & 8.000 & 0.198\\
\bottomrule
\end{tabular}
\end{table}

\endgroup

\section{Cleaning and outlier audit}\label{cleaning-and-outlier-audit}

\tightinputtable{../output/tables/tab_p2_outlier_audit.tex}

\begingroup\scriptsize\setlength{\tabcolsep}{2.5pt}
\begin{table}

\caption{\label{tab:p2-transformation-plan}Training-only transformation and winsorization plan.}
\centering
\begin{tabular}[t]{lrrrrl}
\toprule
Variable & Raw\_Skewness & IQR\_Flags & Winsor\_Lower & Winsor\_Upper & Transformation\\
\midrule
fixed\_acidity & 0.9245 & 31 & 5.1000 & 13.2130 & standardize\\
volatile\_acidity & 0.8077 & 18 & 0.2000 & 1.0400 & standardize\\
citric\_acid & 0.2631 & 0 & 0.0000 & 0.7113 & standardize\\
residual\_sugar & 4.6767 & 98 & 1.4000 & 8.3000 & log1p + standardize\\
chlorides & 5.4213 & 80 & 0.0419 & 0.3732 & log1p + standardize\\
\addlinespace
free\_sulfur\_dioxide & 1.1919 & 21 & 3.0000 & 48.2600 & log1p + standardize\\
total\_sulfur\_dioxide & 1.6105 & 34 & 8.0000 & 144.1300 & log1p + standardize\\
density & 0.0134 & 28 & 0.9918 & 1.0014 & standardize\\
ph & 0.2582 & 23 & 2.9300 & 3.6913 & standardize\\
sulphates & 2.3385 & 46 & 0.4287 & 1.3126 & log1p + standardize\\
\addlinespace
alcohol & 0.8544 & 7 & 9.0000 & 13.5710 & standardize\\
\bottomrule
\end{tabular}
\end{table}

\endgroup

\includefigure{../output/figures/fig_p2_outlier_audit.pdf}{Training observations flagged by the univariate Tukey-IQR rule. Flags are diagnosed rather than deleted wholesale.}{fig:appendix-outliers}

\section{Supplementary selection and model
tables}\label{supplementary-selection-and-model-tables}

\begingroup\scriptsize\setlength{\tabcolsep}{2.5pt}
\begin{table}

\caption{\label{tab:p3-feature-screening}Training-only feature screening and OLS coefficients on the standardized scale.}
\centering
\begin{tabular}[t]{llrrrl}
\toprule
  & Feature & Correlation\_with\_quality & VIF & OLS\_standardized\_coefficient & Stepwise\_AIC\_retained\\
\midrule
11 & alcohol & 0.501 & 3.559 & 0.297 & Yes\\
2 & volatile\_acidity & -0.398 & 1.750 & -0.187 & Yes\\
10 & sulphates & 0.308 & 1.469 & 0.186 & Yes\\
3 & citric\_acid & 0.249 & 3.046 & -0.032 & No\\
7 & total\_sulfur\_dioxide & -0.179 & 3.378 & -0.100 & Yes\\
\addlinespace
8 & density & -0.173 & 7.011 & -0.052 & No\\
1 & fixed\_acidity & 0.136 & 8.075 & 0.054 & No\\
5 & chlorides & -0.134 & 1.435 & -0.087 & Yes\\
4 & residual\_sugar & 0.070 & 1.794 & 0.045 & No\\
9 & ph & -0.070 & 3.423 & -0.064 & Yes\\
\addlinespace
6 & free\_sulfur\_dioxide & -0.067 & 3.013 & 0.052 & Yes\\
\bottomrule
\end{tabular}
\end{table}

\endgroup

\begingroup\scriptsize\setlength{\tabcolsep}{2.5pt}
\begin{table}

\caption{\label{tab:p4-coefficient-comparison}Coefficients at minimum-CV penalties on the standardized transformed scale.}
\centering
\begin{tabular}[t]{lrrrr}
\toprule
Feature & OLS & Ridge & Lasso & Elastic\_Net\\
\midrule
fixed\_acidity & 0.0537 & 0.0594 & 0.0000 & 0.0000\\
volatile\_acidity & -0.1866 & -0.1763 & -0.1773 & -0.1770\\
citric\_acid & -0.0322 & -0.0134 & 0.0000 & 0.0000\\
residual\_sugar & 0.0448 & 0.0466 & 0.0142 & 0.0149\\
chlorides & -0.0873 & -0.0845 & -0.0771 & -0.0777\\
\addlinespace
free\_sulfur\_dioxide & 0.0517 & 0.0405 & 0.0060 & 0.0086\\
total\_sulfur\_dioxide & -0.1004 & -0.0906 & -0.0558 & -0.0586\\
density & -0.0525 & -0.0670 & 0.0000 & 0.0000\\
ph & -0.0639 & -0.0493 & -0.0604 & -0.0609\\
sulphates & 0.1863 & 0.1795 & 0.1650 & 0.1653\\
\addlinespace
alcohol & 0.2970 & 0.2759 & 0.3235 & 0.3222\\
\bottomrule
\end{tabular}
\end{table}

\endgroup

\tightinputtable{../output/tables/tab_p4_enet_alpha_search.tex}

\tightinputtable{../output/tables/tab_p5_error_by_quality.tex}

\section{Reproducibility and
collaboration}\label{reproducibility-and-collaboration}

The canonical analysis is rebuilt from RStudio with
\texttt{source("analysis/00\_run\_all.R")}, or from a terminal with
\texttt{Rscript\ -\/-vanilla\ analysis/00\_run\_all.R}. The five
workstreams execute in dependency order, every random operation uses a
recorded seed, and \texttt{output/artifact\_manifest.csv} stores file
sizes and MD5 checksums. Row-level holdout outcomes are not serialized.
Software versions for the verified run are listed below; full
\texttt{sessionInfo()} is in \texttt{output/logs/session\_info.txt}. The
analysis and dynamic report use R, \texttt{knitr}, and
\texttt{rmarkdown} (\citeproc{ref-RCore2026}{R Core Team 2026};
\citeproc{ref-XieKnitr2025}{Xie 2025};
\citeproc{ref-AllaireRmarkdown2026}{{Allaire et al.} 2026}).

\tightinputtable{../output/tables/tab_reproducibility_versions.tex}

The workstream allocation and review interfaces are documented in the
repository's collaboration work-split file. These organizational labels
must be checked against the members' actual contributions before final
submission.

\newpage

\section*{References}\label{references}
\addcontentsline{toc}{section}{References}

\protect\phantomsection\label{refs}
\begin{CSLReferences}{1}{1}
\bibitem[\citeproctext]{ref-AllaireRmarkdown2026}
{Allaire, JJ, Yihui Xie, Christophe Dervieux, et al.} 2026.
\emph{Rmarkdown: Dynamic Documents for r}.
\url{https://github.com/rstudio/rmarkdown}.

\bibitem[\citeproctext]{ref-Cortez2009}
Cortez, Paulo, António Cerdeira, Fernando Almeida, Telmo Matos, and José
Reis. 2009a. {``Modeling Wine Preferences by Data Mining from
Physicochemical Properties.''} \emph{Decision Support Systems} 47 (4):
547--53. \url{https://doi.org/10.1016/j.dss.2009.05.016}.

\bibitem[\citeproctext]{ref-UCIWine}
Cortez, Paulo, António Cerdeira, Fernando Almeida, Telmo Matos, and José
Reis. 2009b. \emph{Wine Quality}. UCI Machine Learning Repository.
\url{https://archive.ics.uci.edu/dataset/186/wine+quality}.

\bibitem[\citeproctext]{ref-Friedman2010}
Friedman, Jerome, Trevor Hastie, and Robert Tibshirani. 2010.
{``Regularization Paths for Generalized Linear Models via Coordinate
Descent.''} \emph{Journal of Statistical Software} 33 (1): 1--22.
\url{https://doi.org/10.18637/jss.v033.i01}.

\bibitem[\citeproctext]{ref-Hoerl1970}
Hoerl, Arthur E., and Robert W. Kennard. 1970. {``Ridge Regression:
Biased Estimation for Nonorthogonal Problems.''} \emph{Technometrics} 12
(1): 55--67. \url{https://doi.org/10.1080/00401706.1970.10488634}.

\bibitem[\citeproctext]{ref-RCore2026}
R Core Team. 2026. \emph{R: A Language and Environment for Statistical
Computing}. R Foundation for Statistical Computing.
\url{https://doi.org/10.32614/R.manuals}.

\bibitem[\citeproctext]{ref-Tibshirani1996}
Tibshirani, Robert. 1996. {``Regression Shrinkage and Selection via the
Lasso.''} \emph{Journal of the Royal Statistical Society: Series B} 58
(1): 267--88. \url{https://doi.org/10.1111/j.2517-6161.1996.tb02080.x}.

\bibitem[\citeproctext]{ref-XieKnitr2025}
Xie, Yihui. 2025. \emph{Knitr: A General-Purpose Package for Dynamic
Report Generation in r}. \url{https://yihui.org/knitr/}.

\bibitem[\citeproctext]{ref-Zou2005}
Zou, Hui, and Trevor Hastie. 2005. {``Regularization and Variable
Selection via the Elastic Net.''} \emph{Journal of the Royal Statistical
Society: Series B} 67 (2): 301--20.
\url{https://doi.org/10.1111/j.1467-9868.2005.00503.x}.

\end{CSLReferences}

\end{document}
